% Einheit 3 von 3 – Punktsymmetrie, Drehung, Verschiebung (bank/symmetrie-abbildungen/e3.jsonl, zusammenbau v0.1)
%% TODO Zweigzeile Teil 1: was hier gelernt wird (Satz in Schülersprache aus dem Einheitstitel) – Einheit 3
\einheitenkopf[e3]{Einheit 3 von 3 · Punktsymmetrie, Drehung, Verschiebung}
\zweigzeile{ab Kl. 5, je nach Buch bis Kl. 8 · keine P10-Aufgabe}
\setcounter{aufgabe}{21}

%% TODO Titel: Ich-kann-Satz fehlt in der Bank (Platzhalter: Kettenname) – Nr. 22
%% TODO Anweisung: ein Satz über der ersten Teilaufgabe fehlt in der Bank – Nr. 22
\begin{aufgabe}{Drehen, Verschieben, Punktsymmetrie}
%% TODO \rechenplatz{n}: Schrittzahl des Grundfalls fehlt in der Bank (nur bei fester Schreibform, 2.3 b)
\begin{teile}
\teil Wie ist das Dreieck A'B'C' aus dem Dreieck ABC entstanden? \\ \kreuz{an einer Geraden gespiegelt} \\ \kreuz{gedreht} \\ \kreuz{verschoben}

\begin{ksys}[xmin=0,xmax=8,ymin=0,ymax=5] \punkt{1}{1}{A} \punkt{3}{1}{B} \punkt{1}{3}{C} \punkt{5}{2}{A'} \punkt{7}{2}{B'} \punkt{5}{4}{C'} \end{ksys}
\teil Der Verschiebungspfeil führt von P nach Q. Verschiebe das Dreieck ABC mit diesem Pfeil.

\begin{ksys}[xmin=0,xmax=10,ymin=0,ymax=8] \punkt{1}{1}{A} \punkt{3}{1}{B} \punkt{2}{3}{C} \punkt{5}{0}{P} \punkt{9}{1}{Q} \end{ksys}
\teil Das Dreieck ABC wurde auf A'B'C' verschoben. Beschreibe die Verschiebung in Kästchen.

\begin{ksys}[xmin=0,xmax=9,ymin=0,ymax=8] \punkt{1}{4}{A} \punkt{3}{4}{B} \punkt{1}{6}{C} \punkt{5}{2}{A'} \punkt{7}{2}{B'} \punkt{5}{4}{C'} \end{ksys}
\leerfeld[Kästchen] nach \leerfeld , \leerfeld[Kästchen] nach \leerfeld
\teil Drehe das Dreieck ABC um den Punkt Z um eine halbe Drehung.

\begin{ksys}[xmin=-5,xmax=8,ymin=-5,ymax=8] \punkt{5}{5}{A} \punkt{7}{5}{B} \punkt{5}{7}{C} \punkt{4}{4}{Z} \end{ksys}
\teil Drehe das Dreieck ABC um den Punkt Z um eine Vierteldrehung gegen den Uhrzeigersinn.

\begin{ksys}[xmin=-6,xmax=7,ymin=-2,ymax=7] \punkt{1}{1}{A} \punkt{4}{1}{B} \punkt{1}{3}{C} \punkt{0}{0}{Z} \end{ksys}
\teil Spiegle den Punkt P am Zentrum Z.

\begin{ksys}[xmin=-5,xmax=8,ymin=-5,ymax=6] \punkt{1}{2}{P} \punkt{4}{3}{Z} \end{ksys}
P'( \leerfeld | \leerfeld )
\teil Spiegle das Dreieck ABC am Zentrum Z.

\begin{ksys}[xmin=-6,xmax=8,ymin=-5,ymax=8] \punkt{1}{5}{A} \punkt{3}{6}{B} \punkt{2}{7}{C} \punkt{4}{4}{Z} \end{ksys}
\teil Ist das Parallelogramm punktsymmetrisch? \\ \kreuz{ja} \\ \kreuz{nein}

\parallelogramm{4}{2}{55}
\teil Die Dreiecke ABC und DEF sind der Anfang eines Bandornaments. Zeichne die nächsten zwei Dreiecke und benenne die Abbildung.

\begin{ksys}[xmin=0,xmax=12,ymin=0,ymax=4] \punkt{0}{1}{A} \punkt{2}{1}{B} \punkt{1}{3}{C} \punkt{3}{1}{D} \punkt{5}{1}{E} \punkt{4}{3}{F} \end{ksys}
\teil Lege mit Steinen wie diesem Parallelogramm ein Parkett aus mindestens sechs Steinen. Mit welcher Abbildung kommst du von einem Stein zum Nachbarn?

\parallelogramm{2}{1.5}{60}
\teil Welche Abbildung führt das Dreieck ABC in A'B'C' über? Begründe, warum AB und A'B' gleich lang sind.

\begin{ksys}[xmin=-5,xmax=5,ymin=-4,ymax=4] \punkt{1}{1}{A} \punkt{4}{1}{B} \punkt{2}{3}{C} \punkt{-1}{-1}{A'} \punkt{-4}{-1}{B'} \punkt{-2}{-3}{C'} \punkt{0}{0}{Z} \end{ksys}
\end{teile}
\end{aufgabe}

%% TODO Titel: Ich-kann-Satz fehlt in der Bank (Platzhalter: Kettenname) – Nr. 23
%% TODO Anweisung: ein Satz über der ersten Teilaufgabe fehlt in der Bank – Nr. 23
\begin{aufgabe}{drehsymmetrische Figuren erkennen und den Drehwinkel angeben (Vorrat)}
\begin{teile}
\teil Um welchen kleinsten Winkel musst du ein gleichseitiges Dreieck um seinen Mittelpunkt drehen, damit es wieder auf sich selbst liegt? \leerfeld °
\end{teile}
\end{aufgabe}

%% TODO Titel: Ich-kann-Satz fehlt in der Bank (Platzhalter: Kettenname) – Nr. 24
%% TODO Anweisung: ein Satz über der ersten Teilaufgabe fehlt in der Bank – Nr. 24
\begin{aufgabe}{Original und Bild vergleichen: Längen und Winkel bleiben gleich, kongruente Figuren}
\begin{teile}
\teil Das Dreieck ABC wird verschoben. Es ist AB = 4{,}2\,cm und der Winkel bei A misst $50^\circ$. Wie lang ist A'B', wie groß der Winkel bei A'? A'B' = \leerfeld[cm] , Winkel bei A' = \leerfeld °
\end{teile}
\end{aufgabe}

%% TODO Titel: Ich-kann-Satz fehlt in der Bank (Platzhalter: Kettenname) – Nr. 25
%% TODO Anweisung: ein Satz über der ersten Teilaufgabe fehlt in der Bank – Nr. 25
\begin{aufgabe}{Drehen, Verschieben, Punktsymmetrie – Fehler finden, Begründen, Darstellungswechsel, Anwendung}
\begin{teile}
\teil Finde den Fehler. Lina soll das Dreieck ABC um Z um eine halbe Drehung drehen. Sie spiegelt es an der senkrechten Geraden durch Z.
\teil Begründe, warum Original und Bild bei einer Verschiebung deckungsgleich sind.
\teil Ein Punkt wird von $A(1|2)$ nach $A'(5|3)$ verschoben. Beschreibe die Verschiebung in Worten.
\teil Ein Riesenrad hat 30\,m Durchmesser. Eine Gondel ist ganz unten, 2\,m über dem Boden. Das Rad dreht sich um eine halbe Drehung. Wie hoch ist die Gondel jetzt? \leerfeld[m]
\end{teile}
\end{aufgabe}
